\section{The initial state and its comparison forms}
\label{sec:initialization}

This section constructs the initial state with every site active and its
partial potential equal to zero, for every fixed $d\ge3$ and the
corresponding interval $I_d$ in \eqref{eq:energy-windows}. We prove that the initial inverse has a
bounded diagonal, a small squared-entry operator, small fourth-power row
sums, and the three comparison bounds in \eqref{eq:comparators}. We first establish the lattice
inequalities used to express those bounds and to localize them later.

\subsection{Duality and lattice estimates}

Let $\mathbb T_L=(\ZZ/L\ZZ)^d$, $L\ge4$, and $N=L^d$. An index set
$\mathcal A$ consists of sites of $\mathbb T_L$, with at most two copies
of each site permitted.  Write $\pi:\mathcal A\to\mathbb T_L$ for the
underlying site map.  For a function $v$ on $\mathbb T_L$, let
$(\mathcal Cv)_x=v_{\pi(x)}$.  The adjoint sends a charge $y$ on
$\mathcal A$ to
\[
 (y_\Sigma)_i=(\mathcal C^*y)_i
 =\sum_{x:\pi(x)=i}y_x,
\]
with value zero at sites not represented in $\mathcal A$.
The positive nearest-neighbor Laplacian is denoted by $\Lambda$, and
$\Lambda^+$ is its inverse on mean-zero functions, extended by zero on
constants.  We normalize the gradient energy by
$\|\nabla v\|_2^2=\langle v,\Lambda v\rangle$.

\begin{lemma}[Dual comparison forms]\label{lem:dual-forms}
For $b>0$, define
\[
 \begin{split}
 W(f)&=\inf_{f=\alpha+\mathcal Cv}
       \bigl(\|\alpha\|_2^2+b^{-1}\|\nabla v\|_2^2\bigr),\\
 W^*(y)&=\|y\|_2^2+
       b\langle y_\Sigma,\Lambda^+y_\Sigma\rangle,
       \qquad \sum_{x\in\mathcal A}y_x=0.
 \end{split}
\]
Then $W$ is a positive definite quadratic form on functions on
$\mathcal A$ modulo constants, and $W^*$ is its dual quadratic form:
\[
 W^*(y)=\sup_f\bigl(2\Re\langle y,f\rangle-W(f)\bigr).
\]
In particular, if $D$ is a Laplacian on $\mathcal A$ whose kernel is
exactly the constants and $X=D^+$, then
$cW^*\le X\le CW^*$ on neutral charges is equivalent to
$C^{-1}W\le D\le c^{-1}W$ on functions modulo constants.
\end{lemma}

\begin{proof}
Because $\mathcal C$ carries constants to constants, $W$ is unchanged by
addition of a constant to $f$.  The minimization is a finite-dimensional
least-squares problem, so its minimum is attained and is a quadratic
form in $f$.  Zero energy forces $\alpha=0$ and $v$
constant, so the induced form on the quotient is positive definite.

For a neutral $y$, completion of squares gives
\[
\begin{split}
 \sup_f\bigl(2\Re\langle y,f\rangle-W(f)\bigr)
 &=\sup_{\alpha,v}\bigl(
       2\Re\langle y,\alpha+\mathcal Cv\rangle
       -\|\alpha\|_2^2-b^{-1}\langle v,\Lambda v\rangle\bigr)\\
 &=\|y\|_2^2+b\langle\mathcal C^*y,
                         \Lambda^+\mathcal C^*y\rangle.
\end{split}
\]
Neutrality makes $\mathcal C^*y$ orthogonal to constants.  The inverse
matrix of a positive quadratic form is its dual, which proves both the
identity and the comparison assertion on the quotient.
\end{proof}

We use unnormalized Fourier coefficients
$\widehat y(k)=\sum_x y_xe^{-ik\cdot x}$, with inverse normalization
$N^{-1}$.  Frequencies belong to the grid
$\mathcal K_L=(2\pi/L)\ZZ^d$ modulo $2\pi\ZZ^d$, represented in
$[-\pi,\pi)^d$; the symbol of
$\Lambda$ is
\begin{equation}\label{eq:initial-laplacian-symbol}
 \ell(k)=4\sum_{j=1}^d\sin^2(k_j/2)\asymp |k|^2.
\end{equation}

\begin{lemma}[A charge estimate for fixed dimension $d\ge3$]
\label{lem:fourier-charge}
There is a constant $C_d$ such that, for every $L\ge4$, every index
set $\mathcal A$ as above, and every charge $y$ on $\mathcal A$,
\[
 \langle y_\Sigma,\Lambda^+y_\Sigma\rangle
 \le C_d\|y\|_1^{4/3}\|y\|_2^{2/3}.
\]
\end{lemma}

\begin{proof}
Counting grid frequencies in dyadic annuli and using
\eqref{eq:initial-laplacian-symbol} gives
\begin{equation}\label{eq:initial-frequency-sum}
 N^{-1}\sum_{0<|k|\le a}\ell(k)^{-1}
 \le C_d a^{d-2}\le C_d a,
 \qquad 0<a\le1.
\end{equation}
For completeness, a ball of radius $r\ge2\pi/L$ contains at most
$C_d(1+Lr)^d\le C_dNr^d$ grid points. On an annulus
$r/2<|k|\le r$, the summand is at most $C r^{-2}$. Decomposing the
ball of radius $a$ into dyadic annuli, discarding empty annuli below
the first mesh frequency, gives
$C_d\sum_{j\ge0}(2^{-j}a)^{d-2}\le C_da^{d-2}$.
If $a$ is smaller than the first nonzero frequency, the sum is empty.
The weaker last bound deliberately retains the exponents needed later;
it is valid because $a\le1$ and $d-2\ge1$.

At each frequency $|\widehat{y_\Sigma}(k)|\le\|y\|_1$, while at most
two copies per site gives $\|y_\Sigma\|_2^2\le2\|y\|_2^2$.
The low frequencies in the quadratic form are therefore bounded by
$Ca\|y\|_1^2$, and Parseval bounds the high frequencies by
$Ca^{-2}\|y\|_2^2$.  For $y\ne0$, choose
$a=(\|y\|_2/\|y\|_1)^{2/3}\le1$.  This proves the claimed
inequality.  For $y=0$ there is nothing to prove.
\end{proof}

\begin{lemma}[Weighted cube inequality]\label{lem:cube-weight}
Let $Q$ be a lattice cube in $\ZZ^d$ of integer side $l\ge1$,
let $v:Q\to\CC$, and let $w_i\ge0$ be weights. With the gradient
restricted to edges contained in $Q$,
\[
 \|v-\operatorname{avg}_Qv\|_{2,Q}^2
 \le C_d l^2\|\nabla v\|_{2,Q}^2,
\]
and
\begin{equation}\label{eq:initial-weighted-cube}
 \sum_{i\in Q}w_i|v_i-\operatorname{avg}_Qv|^2
 \le C_d(\max_Q w)^{1/3}\Bigl(\sum_Qw\Bigr)^{2/3}
          \|\nabla v\|_{2,Q}^2.
\end{equation}
For fixed $d$, the constants are independent of $l$, $v$, and the
weights. The statements also hold for cubes in a torus whenever their
internal edges are identified with those of a lattice cube.
\end{lemma}

\begin{proof}
The case $l=1$ is immediate. Translate $Q$ to
$\{0,\ldots,l-1\}^d$. Define
\[
 \rho(j)=
 \begin{cases}j,&0\le j<l,\\2l-1-j,&l\le j<2l,\end{cases}
 \qquad
 V(j_1,\ldots,j_d)=v(\rho(j_1),\ldots,\rho(j_d))
                         -\operatorname{avg}_Qv.
\]
This gives a centered function on $\mathbb T_{2l}$. Each original site
has $2^d$ preimages. Every internal cube edge is copied $2^d$ times,
and differences on the folding edges vanish. Consequently
\[
 \|V\|_2^2=2^d\|v-\operatorname{avg}_Qv\|_{2,Q}^2,
 \qquad \Lambda_{2l}(V)=2^d\|\nabla v\|_{2,Q}^2.
\]
The smallest positive eigenvalue of $\Lambda_{2l}$ is at least
$c l^{-2}$, so Parseval proves the first assertion.

For the weighted assertion put $M=\max_Q w$ and $S=\sum_Qw$.
If $S=0$ there is nothing to prove. Otherwise
$a=(M/S)^{1/3}\in(0,1]$. Split $V$ into frequencies $0<|k|\le a$
and $|k|>a$; its zero Fourier coefficient vanishes. With normalized
inverse Fourier transformation, Cauchy--Schwarz gives
\[
 |V_{\rm low}(x)|^2
 \le\left((2l)^{-d}\sum_{0<|k|\le a}\ell(k)^{-1}\right)
       \left((2l)^{-d}\sum_k\ell(k)|\widehat V(k)|^2\right).
\]
Thus \eqref{eq:initial-frequency-sum} and Parseval imply
\[
 \|V_{\rm low}\|_\infty^2\le C_d a\Lambda_{2l}(V),
 \qquad
 \|V_{\rm high}\|_2^2\le C_d a^{-2}\Lambda_{2l}(V).
\]
These inequalities include the case of an empty low-frequency range.
Restricting to $Q$, using $|z_1+z_2|^2\le2|z_1|^2+2|z_2|^2$,
and applying the exact reflection identities above, we obtain
\[
 \sum_{i\in Q}w_i|v_i-\operatorname{avg}_Qv|^2
 \le C_d(Sa+Ma^{-2})\|\nabla v\|_{2,Q}^2.
\]
The chosen value of $a$ gives
$Sa=Ma^{-2}=M^{1/3}S^{2/3}$, proving the claim.
\end{proof}

\begin{lemma}[Spatial counting and exponential moments]
\label{lem:lattice-counting}
For ordinary nearest-neighbor distance on $\mathbb T_L$, balls satisfy
$|B(x,R)|\le C_d(1+R)^d$. If $1\le R\le L/4$, they also satisfy
$|B(x,R)|\ge c_dR^d$. Uniformly in $L$, for $0<u\le1$ and $a=0,2$,
\begin{equation}\label{eq:lattice-exponential-moments}
 \sum_y\dist(x,y)^a e^{-u\dist(x,y)}\le C_d u^{-d-a},
\end{equation}
where the factor $\dist(x,x)^0$ is interpreted as $1$.
\end{lemma}

\begin{proof}
Represent each displacement by coordinate differences of least
absolute value. A distance ball is contained in the coordinate box
$[-R,R]^d$, proving the upper bound. For $R\le L/4$, a coordinate
box of radius $\lfloor R/d\rfloor$ embeds without identifications
and is contained in the distance ball, giving the lower bound after
adjusting the constant for bounded $R$. The number of representatives
at distance $r$ is at most $C_d(1+r)^{d-1}$, by counting nonnegative
integer $d$-tuples of sum $r$ and their signs. Thus the left side of
\eqref{eq:lattice-exponential-moments} is bounded by
$C_d\sum_{r\ge0}(1+r)^{d-1+a}e^{-ur}$, which is at most
$C_du^{-d-a}$ by comparison with its integral.
\end{proof}

These lattice estimates are uniform in the torus and cube sizes.
They will allow local changes to be measured in the same forms as the
initial state.  We now construct that state and identify its Fourier
comparison operator.

\subsection{The fixed interval and the scalar diagonal equation}

Self-consistent diagonal equations and free lattice resolvent estimates
also play a central role in
\cite[Section~2.2 and Appendix~A]{BlackDroginHernandez2025}.
We prove the specific finite-volume estimates needed here, including
the gap at the smallest nonzero grid frequency.

The free dispersion in the negative-hopping convention is
\[
 a(k)=-2\sum_{j=1}^d\cos k_j=-2d+\ell(k).
\]
Write $a_d=e_2$ for $d=3$ and $a_d=1/100$ for $d\ge4$, so that
$I_d=-2d+(a_d/2,a_d)$. Choose $e_2>0$ fixed and sufficiently small,
with $e_2<1/100$. For estimates involving a small shift of energy,
use the larger compact offset interval
\[
 \mathcal J_d=[2a_d/5,5a_d/4]\subset(0,0.04).
\]
For $d\ge4$ this is exactly $[1/250,1/80]$.
The parameter $a_d$ is fixed before the coupling or any scale is chosen.
In the free-resolvent analysis below, constants may depend on $d$ and this
fixed $a_d$; this dependence is included in the notation $C_d$.
In what follows, integrals over the frequency torus are
normalized by $(2\pi)^{-d}$.

\begin{lemma}[The free resolvent near the band minimum]
\label{lem:initial-bottom-density}
The function
\[
 G_0(z)=\int_{[-\pi,\pi)^d}\frac{1}{a(k)-z}\,
                    \frac{dk}{(2\pi)^d},\qquad \Im z>0,
\]
extends continuously from above to $-2d+\mathcal J_d$. Its boundary imaginary
part there is bounded above and below by positive constants depending
only on the fixed dimension and window. Uniformly for
$\alpha\in-2d+\mathcal J_d$ and $0<v\le1$,
\begin{equation}\label{eq:initial-density-integrals}
 \int\frac{dk/(2\pi)^d}{|a(k)-\alpha-iv|}
       \le C_d\log(2/v),
 \qquad
 \left|\int\frac{dk/(2\pi)^d}{(a(k)-\alpha-iv)^2}\right|
       \le C_d\log(2/v).
\end{equation}
The second estimate holds uniformly also for
$-2\sum_{j=1}^dc_j\cos(k_j+\phi_j)$ whenever
$\max_j|c_j-1|\le\varepsilon_d$ for a sufficiently small fixed
$\varepsilon_d>0$, with arbitrary real phases $\phi_j$.
\end{lemma}

\begin{proof}
On $(-\pi,\pi)^d$ make the one-to-one change of coordinates
$x_j=2\sin(k_j/2)$. Then $a(k)+2d=|x|^2$ and
$dk=\prod_j(1-x_j^2/4)^{-1/2}dx$. For $0<s<0.04$ the entire
sphere of radius $\sqrt s$ lies in the coordinate domain, so polar
integration shows that the normalized density of $a+2d$ equals
\begin{equation}\label{eq:initial-density-formula}
 \rho_d(s)=\frac{s^{d/2-1}}{2(2\pi)^d}
   \int_{S^{d-1}}\prod_{j=1}^d
       (1-s\omega_j^2/4)^{-1/2}\,d\sigma(\omega).
\end{equation}
This density is positive and smooth on $(0,0.04)$. On every compact
subinterval its derivatives are bounded and the density has a positive
lower bound, with constants allowed to depend on $d$.

Choose a smooth cutoff $\chi$ supported in $(a_d/5,2a_d)$ and equal
to one on a neighborhood of $\mathcal J_d$. Split the Stieltjes integral into
this localized density and the complementary part. In the latter
part the real denominator stays a fixed distance from zero for
energies in $-2d+\mathcal J_d$, so it is continuous and bounded up to the real
axis, as is its derivative. In the former part, subtract
$\rho_d(\alpha+2d)$ locally from the density. The difference vanishes
linearly at the real energy; its Cauchy integral has a continuous
upper boundary value. The constant part is an explicit logarithmic
integral. This proves upper-half-plane continuity and gives the
boundary imaginary part $\pi\rho_d(\alpha+2d)$.

For the first inequality in \eqref{eq:initial-density-integrals},
integrate the bounded localized density against
$((s-(\alpha+2d))^2+v^2)^{-1/2}$; its integral over the fixed support
is at most $C\log(2/v)$. For the second, integration by parts gives
\[
 \int\frac{\chi(s)\rho_d(s)}{(s-(\alpha+2d)-iv)^2}\,ds
 =\int\frac{(\chi\rho_d)'(s)}{s-(\alpha+2d)-iv}\,ds.
\]
There is no boundary term. The same absolute-integral estimate bounds
this expression, while the complementary contribution stays bounded.

For the perturbation assertion, first translate all phases in the
continuum torus integral. Put
$x_j=2\sqrt{c_j}\sin(k_j/2)$, so the new minimum is
$-2\sum_jc_j$ and the excess energy is $|x|^2$. The Jacobian is
\[
 \prod_{j=1}^d c_j^{-1/2}
            (1-x_j^2/(4c_j))^{-1/2}.
\]
Choose $\varepsilon_d$ so small that $c_j\ge1/2$ and
$2d\varepsilon_d<a_d/10$. Energies in $-2d+\mathcal J_d$ then have excess
above the shifted minimum in a fixed compact subinterval of
$(a_d/5,2a_d)$. The corresponding polar density and its derivatives
are uniformly bounded there. A cutoff equal to one on that compact
interval gives the same integration-by-parts estimate, uniformly in
the coefficients and phases.
\end{proof}

\begin{proposition}[Initialization]\label{prop:initialization}
For each fixed $d\ge3$ there are $\lambda_0(d)>0$ and positive
constants $a_0,B,c,C,C_0$, depending only on the fixed dimension and
window, with the following
property.  For
$0<\lambda<\lambda_0(d)$, the scales in \eqref{eq:scales},
$m\ge4n_0$, $L=4^{20m}$, and every $E\in I_d$,
write $t=t_{n_0}$.  There is a scalar $r$ with
\[
 |\Re r|\le B,\qquad a_0\le p:=\Im r\le B,
 \qquad
 r=N^{-1}\sum_k\frac1{a(k)-E-tr}.
\]
It is the unique solution in the displayed rectangle and depends
continuously on $E$.

Take all sites active, all partial potentials zero, and
$g=(H_0-E-tr)^{-1}$, where $H_0$ is the free hopping matrix on
$\mathbb T_L$.  With $h,K,F,D,X$ as in
Section~\ref{sec:state}, one has $g_{ii}=r$, $\|K\|\le1/10$, and
\[
 \max_i\sum_j|h_{ij}|^4\le\tfrac1{10}t^{1+\delta}.
\]
The kernel of $D$ is exactly the constants, and on neutral charges,
\begin{equation}\label{eq:initial-comparisons}
 \begin{split}
 X&\ge ct^2\Lambda^+,\qquad X\ge c\Id,\\
 X&\le C_0\Id+Ct^2\Lambda^+,\\
 X&\le(1+t^\kappa)p^{-2}\Id+Ct\Lambda^+.
 \end{split}
\end{equation}
Thus this state satisfies \eqref{eq:comparators}, with startup
$b=t_{n_0}t=t^2$ and constants independent of $E,\lambda,m$.
\end{proposition}

\begin{proof}
We first solve the scalar equation, then prove the two smallness bounds,
and finally identify the inverse form of $D$.

\paragraph{The scalar fixed point.}
By Lemma~\ref{lem:initial-bottom-density}, the boundary values
$G_0(\alpha+i0)$ for $\alpha\in-2d+\mathcal J_d$ form a compact subset
of the upper half-plane. Choose $a_0>0$ and $B<\infty$ so that they lie
strictly inside the rectangle
\[
 \mathcal B=\{r:|\Re r|\le B,\ a_0\le\Im r\le B\}.
\]
The closure of $I_d$ lies in the interior of $-2d+\mathcal J_d$.
Hence, when $t$ is small, $E+t\Re r\in-2d+\mathcal J_d$ for every
$E\in I_d$ and $r\in\mathcal B$. Uniform boundary continuity on this
compact interval shows that $r\mapsto G_0(E+tr)$ maps $\mathcal B$
into its interior with fixed slack.

For a periodic $C^1$ function, comparing each grid cell with its corner
gives an integral-to-grid error at most $CL^{-1}\|\nabla f\|_\infty$.
On $\mathcal B$, set $z=E+tr$ and $R(k)=(a(k)-z)^{-1}$.
Since $\Im z\asymp t$,
\begin{equation}\label{eq:initial-grid-derivatives}
 \|\nabla R\|_\infty\le Ct^{-2},\qquad
 \|\nabla(R^2)\|_\infty\le Ct^{-3}.
\end{equation}
The finite-grid map
$\Phi(r)=N^{-1}\sum_kR(k)$ consequently preserves $\mathcal B$,
provided $C/(Lt^2)$ is sufficiently small.  Here
$t\asymp4^{-n_0}$ and $L\ge4^{20n_0}$, so
\begin{equation}\label{eq:initial-grid-error}
 \varepsilon_L:=\frac{C}{Lt^2}\le Ct^{18}.
\end{equation}

The second estimate of Lemma~\ref{lem:initial-bottom-density} holds for
every $r\in\mathcal B$, without a fixed-point hypothesis.  Differentiating
the finite sum and using \eqref{eq:initial-grid-derivatives} gives
\[
 \sup_{r\in\mathcal B}|\Phi'(r)|
 \le Ct\log(1/t)+\varepsilon_L.
\]
It is less than $1/2$ for small $t$.  Integrating the complex derivative
along any segment in the convex rectangle makes $\Phi$ a contraction.
Iteration from one fixed starting point in $\mathcal B$ constructs its unique
fixed point.  The iterates are continuous in $E$ and converge uniformly,
so the fixed point is continuous in $E$.

At this solution, translation invariance gives $g_{ii}=r$.  Taking
imaginary parts of the scalar equation, and cancelling $p>0$, yields
the exact finite-grid normalization
\begin{equation}\label{eq:initial-normalization}
 t\operatorname{avg}_k|R(k)|^2=1.
\end{equation}

\paragraph{The Fourier symbols and the remaining bounds.}
For each grid displacement $\zeta\in\mathcal K_L$, Fourier transformation gives
\begin{equation}\label{eq:initial-symbols}
 \widehat F(\zeta)
    =t\operatorname{avg}_k R(k)\overline{R(k+\zeta)},\qquad
 \widehat K(\zeta)
    =t\operatorname{avg}_k R(k)R(k+\zeta).
\end{equation}
The first symbol is real, either because $F$ is real and even or by
the change $k\mapsto-k-\zeta$ and $R(-k)=R(k)$. Its zero mode is
$1$ by \eqref{eq:initial-normalization}, and Cauchy--Schwarz gives
$|\widehat F(\zeta)|\le1$. Both convolution operators are unitarily
diagonalized by Fourier transformation, and Parseval gives the row identity:
\[
 \|K\|=\max_{\zeta\in\mathcal K_L}|\widehat K(\zeta)|,\qquad
 \sum_j|h_{ij}|^4
   =N^{-1}\sum_{\zeta\in\mathcal K_L}|\widehat F(\zeta)|^2.
\]
Thus the two smallness assertions require a uniform bound for
$\widehat K$ and an averaged square bound for $\widehat F$.
At the fixed point $q=p$ is scalar, so $D=p^2(\Id-F)$.
The inverse-form comparisons will follow from a two-sided estimate of
$1-\widehat F(\zeta)$ at every nonzero grid frequency. At higher
frequencies we will also use that $\widehat F$ is small to retain the
sharp diagonal coefficient.

\paragraph{An absolute-product bound, including tangential intersections.}
We first pass to the continuum frequency torus. For every nonzero
displacement $\zeta\in[-\pi,\pi)^d$, we claim
\begin{equation}\label{eq:initial-absolute-product}
 t\int |R(k)R(k+\zeta)|\,\frac{dk}{(2\pi)^d}
 \le\frac{C_dt\log^2(1/t)}{|\zeta|}.
\end{equation}
For quadratic dispersion, Erd\H{o}s, Salmhofer and Yau obtain a related
inverse-displacement two-resolvent bound by radial and angular integration
\cite[author-preprint Appendix~A, Lemma~A.1, equations~(A.2) and~(A.4)]{ErdosSalmhoferYau2008}.
The proof below supplies the required near-bottom lattice estimate by
convex slicing, including tangency.

In this estimate the displacement need not lie on the finite grid.
For small $t$ the real excess energy $\Re z+2d$ lies in $\mathcal J_d$.
If either dispersion has excess at least $0.04$, the corresponding
denominator is separated from zero by a fixed positive constant.
Lemma~\ref{lem:initial-bottom-density} then bounds this part of the
integral, including the factor $t$, by $C_dt\log(1/t)$. This is
absorbed by the right side of \eqref{eq:initial-absolute-product},
since $|\zeta|\le\pi\sqrt d$.

In the joint region where both excess energies are below $0.04$,
the frequency vectors $k$ and $k+\zeta$ have unique small unwrapped
representatives: every coordinate has absolute value at most
$2\arcsin(0.1)<0.201$. Their difference is the representative
$\zeta$, so $|\zeta_j|<0.402$. Set
$x_j=\sin(k_j+\zeta_j/2)$. The joint region is contained in
$|x_j|<1/4$, where the change of variables and its inverse have bounded
Jacobians. The two dispersions become
\[
 A_\pm(x)=-2\sum_{j=1}^d\cos(\zeta_j/2)\sqrt{1-x_j^2}
             \ \pm\ 2\sum_{j=1}^d\sin(\zeta_j/2)x_j.
\]
Since $\cos(\zeta_j/2)\ge\cos(0.201)>0$, their Hessians are
uniformly positive and bounded on that box. Their
difference is linear:
\[
 A_+(x)-A_-(x)=b_\zeta\cdot x,
 \qquad (b_\zeta)_j=4\sin(\zeta_j/2),
 \qquad c|\zeta|\le|b_\zeta|\le C|\zeta|.
\]

Suppose that the two energy constraints
\[
 |A_-(x)-\Re z|\le s_1,\qquad |A_+(x)-\Re z|\le s_2
\]
hold simultaneously, where $s_1,s_2>0$. Their difference gives
$|b_\zeta\cdot x|\le s_1+s_2$. Thus $x$ lies in a slab normal to
$b_\zeta$ of width at most $C\max(s_1,s_2)/|\zeta|$. On each slice
of this slab, the remaining task is to bound the narrower energy shell,
including when the slice is tangent to an energy surface.

For either sign, every affine hyperplane $H$ meeting the box, every
$E'\in\RR$, and every $s>0$, we claim
\[
 \operatorname{vol}_{H}\bigl\{x\in H\cap[-1/4,1/4]^d:
             |A_\pm(x)-E'|\le s\bigr\}\le C_ds,
\]
where $\operatorname{vol}_{H}$ is $(d-1)$-dimensional volume on $H$.
Fix the sign and extend each one-variable summand outside its
coordinate interval by its second-order Taylor quadratic at the
endpoint. This gives a globally $C^2$ function $f$ with
$c\Id\le\nabla^2f\le C\Id$, uniformly in $\zeta$. On a hyperplane
$H$ that meets the box choose $x_0\in H$ in the box. The restriction
$f|_H$ is coercive and uniformly convex, and therefore has a unique
minimizer $x_*$. If $P_H$ is projection to the direction space of $H$,
strong convexity and $P_H\nabla f(x_*)=0$ give
\[
 c|x_*-x_0|\le |P_H\nabla f(x_0)|\le C_d.
\]
Consequently $H$ intersected with the box lies in a ball of a fixed
radius $R_d$ about $x_*$. In polar coordinates in $H$, of dimension
$q=d-1\ge2$, one has
\[
 \frac{d}{dr}f(x_*+r\omega)\ge cr,
 \qquad
 \frac{r^{q-1}}{\frac{d}{dr}f(x_*+r\omega)}
       \le C r^{q-2}\le C_d\quad(0<r\le R_d).
\]
Along each ray the energy is increasing. Changing from $r$ to the
energy and integrating an interval of length $2s$ proves the asserted
shell bound, also when the interval contains the minimum. This
explicitly includes tangency of the hyperplane to an energy surface;
no nonzero-gradient assumption at the constrained minimum is used.

On each slice of the slab above, apply this shell bound to the narrower
of the two energy constraints. Fubini's theorem
and the bounded coordinate Jacobian give joint volume at most
\[
 C_d\frac{\max(s_1,s_2)\min(s_1,s_2)}{|\zeta|}
       =C_d\frac{s_1s_2}{|\zeta|}.
\]
Partition each denominator into dyadic sizes $s_i\ge ct$, with the
innermost bin of size comparable to $t$ and the outermost of fixed
size. On a pair of bins the product inverse is at most
$C/(s_1s_2)$, while the containing shells have the volume just bounded.
There are $O(\log(1/t))$ bins for each denominator. Multiplication
by $t$ and summation prove \eqref{eq:initial-absolute-product}.

Uniformly in $\zeta$,
\[
 \|\nabla(RR_\zeta)\|_\infty+
 \|\nabla(R\overline{R_\zeta})\|_\infty\le C_dt^{-3},
 \qquad R_\zeta(k)=R(k+\zeta).
\]
Thus the error in either $t$-weighted product average is at most
$\varepsilon_L$ from \eqref{eq:initial-grid-error}. These absolute
quadrature errors will be used only for product estimates, not to
prove the gap at very small nonzero frequencies.

\paragraph{Smallness of $K$ and the fourth-power rows.}
The absolute-product estimate is small for $|\zeta|>\sqrt t$.
To control $\widehat K$ also at smaller grid frequencies, use cancellation
in the signed product. For $\zeta\in\mathcal K_L$ with
$|\zeta|\le\sqrt t$, use the elementary identity
\[
 \frac1{AB}=\int_0^1\frac{du}{[(1-u)A+uB]^2}.
\]
The two denominators have the same strictly negative imaginary part,
so none of the interpolated denominators vanishes.  The interpolated
dispersion $(1-u)a(k)+ua(k+\zeta)$ is a cosine sum with shifted phases
and amplitudes satisfying
\[
 A_j(u)^2=1-4u(1-u)\sin^2(\zeta_j/2)=1-O(t).
\]
Lemma~\ref{lem:initial-bottom-density} bounds the integral of its
inverse square by $C\log(1/t)$, uniformly in $u$ and $\zeta$.
Consequently
\[
 |\widehat K(\zeta)|\le Ct\log(1/t)+\varepsilon_L
       \quad (|\zeta|\le\sqrt t).
\]
For larger $\zeta$, \eqref{eq:initial-absolute-product} bounds it by
$C\sqrt t\log^2(1/t)+\varepsilon_L$.  These estimates tend to zero
uniformly, giving the asserted bound on $\|K\|$.

For the fourth-power row sums, the Parseval identity above and the
absolute-product bound give the following estimate. All frequency sums
in this calculation are over $\mathcal K_L$:
\begin{equation}\label{eq:initial-fourth-row}
 \begin{split}
 \sum_j|h_{ij}|^4
 &\le N^{-1}+Ct^2\log^4(1/t)
       \left(N^{-1}\sum_{\zeta\ne0}|\zeta|^{-2}\right)
       +C\varepsilon_L^2\\
 &\le N^{-1}+Ct^2\log^4(1/t)+C\varepsilon_L^2.
 \end{split}
\end{equation}
The frequency sum is bounded uniformly by the $d$-dimensional shell
count used earlier.  Here $N^{-1}\le C_dt^{20d}$ and
$\varepsilon_L^2\le Ct^{36}$; hence the last expression is
$o(t^{1+\delta})$, providing the stated fourth-row slack.

We have obtained a well-defined initial diagonal and the two smallness
bounds.  It remains to prove the gap of $F$ and identify the resulting
inverse form.  The gap at the smallest frequencies requires a direct
estimate of the finite sum, which we give next.

\paragraph{The gap at every grid frequency.}
In this part $\zeta\in\mathcal K_L$.
The exact normalization and \eqref{eq:initial-symbols} imply
\begin{equation}\label{eq:initial-gap-identity}
 1-\widehat F(\zeta)
   =\frac t2\operatorname{avg}_k|R(k)-R(k+\zeta)|^2.
\end{equation}
The dispersion is Lipschitz, so the resolvent identity and
$|R|\le C/t$ give
\[
 |R(k)-R(k+\zeta)|
 \le C|\zeta|\,|R(k)|\,|R(k+\zeta)|.
\]
Bound one factor by $C/t$ and use
\eqref{eq:initial-normalization} for the other.  Since also
$|\widehat F|\le1$, this proves
\begin{equation}\label{eq:initial-gap-upper}
 1-\widehat F(\zeta)\le C\min(1,|\zeta|^2/t^2).
\end{equation}

For the lower bound choose a sufficiently small fixed number
$\zeta_0>0$. Suppose $0<|\zeta|\le\zeta_0$, and put
$e=\zeta/|\zeta|$ and $\alpha=\Re z+2d\in\mathcal J_d$.
In the coordinates $x_j=2\sin(k_j/2)$, choose fixed
$0<\theta_1<\theta_0$ and a fixed small $c_0>0$, and define
\[
 \begin{split}
 P_e&=\{k:\ ||x(k)|^2-\alpha|\le c_0t,\quad
                  x(k)\cdot e/|x(k)|\ge\cos\theta_0\},\\
 P_e^-&=\{k:\ ||x(k)|^2-\alpha|\le c_0t/2,\quad
                  x(k)\cdot e/|x(k)|\ge\cos\theta_1\}.
 \end{split}
\]
All these points are in the bottom coordinate neighborhood. The radii
are bounded above by a small constant and below by a positive constant.
Writing $x=\rho\omega$, the derivative identity
\[
 \nabla a(k)\cdot e
 =2\sum_j x_j\sqrt{1-x_j^2/4}\,e_j
 \ge 2\rho\cos\theta_0-\rho^3/2
\]
uses $|\sqrt{1-u}-1|\le u$ and
$\|(x_j^3)_j\|_2\le\rho^3$. Taking $\cos\theta_0\ge3/4$,
and using $\rho^2<0.02$, proves
$\nabla a(k)\cdot e\ge\rho\ge c_1>0$ on $P_e$.
The constants are uniform
in $e,\alpha$ and $t$. Shrink $\zeta_0$ so that the segment from
$k$ to $k+\zeta$ remains in this small neighborhood. Its Hessian
$2\diag(\cos k_j)$ is positive there, so the directional derivative
along the segment stays at least $c_1$.

We next count grid points directly. Polar integration in the $x$
coordinates, with its bounded positive Jacobian, gives
\[
 (2\pi)^{-d}|P_e^-|\ge c_dt.
\]
The angular margin from $\theta_1$ to $\theta_0$ is fixed, while
the radial energy margin is $c_0t/2$. Both the coordinate map and
its inverse have uniformly bounded derivatives. It follows that
$P_e^-$ lies at Euclidean distance at least $c_dt$ from the complement
of $P_e$: a displacement of size $c_dt$ changes the energy by at
most $c_0t/2$ and the angular coordinate by less than its fixed margin.
The mesh-cell diameter is $2\pi\sqrt d/L=o(t)$. Associate to each
point of $P_e^-$ the corner grid point of its mesh cell. For small $t$
that grid point belongs to $P_e$. Each cell has normalized volume
$N^{-1}$, and hence
\begin{equation}\label{eq:initial-grid-patch-count}
 N^{-1}\#(\mathcal K_L\cap P_e)\ge c_dt.
\end{equation}
This argument is uniform in the direction of $\zeta$, including
directions that vary with the grid size.

At each of these points,
\[
 |a(k+\zeta)-a(k)|\ge c|\zeta|,\qquad
 |a(k)-z|\le Ct,\qquad
 |a(k+\zeta)-z|\le C(t+|\zeta|).
\]
Their nonnegative contribution to \eqref{eq:initial-gap-identity}
therefore gives
\begin{equation}\label{eq:initial-direct-grid-gap}
 1-\widehat F(\zeta)
 \ge c\frac{|\zeta|^2}{(t+|\zeta|)^2},
       \qquad 0<|\zeta|\le\zeta_0.
\end{equation}
This argument is valid even for $|\zeta|$ of order $L^{-1}$: no
additive quadrature error has been subtracted from the gap.

For $|\zeta|\ge\zeta_0$, the absolute-product estimate and the grid
error give
\[
 |\widehat F(\zeta)|
 \le Ct\log^2(1/t)/\zeta_0+\varepsilon_L<1/2
\]
when $t$ is small.  Combining this with
\eqref{eq:initial-gap-upper} and \eqref{eq:initial-direct-grid-gap}
proves
\begin{equation}\label{eq:initial-gap-comparison}
 1-\widehat F(\zeta)\asymp
       \min(1,|\zeta|^2/t^2),\qquad \zeta\ne0.
\end{equation}

\paragraph{The three inverse-form bounds.}
Equation~\eqref{eq:initial-gap-comparison} and $D=p^2(\Id-F)$ now show
that $D$ has precisely the constants as kernel. Thus, for
$\zeta\in\mathcal K_L\setminus\{0\}$,
\[
 \widehat X(\zeta)=\frac{p^{-2}}{1-\widehat F(\zeta)}.
\]
The same gap comparison gives
\[
 c\max(1,t^2/|\zeta|^2)
 \le\widehat X(\zeta)
 \le C(1+t^2/|\zeta|^2).
\]
Together with \eqref{eq:initial-laplacian-symbol}, this proves the two
rough comparisons and the lower bound by $c\Id$.

For the sharp upper bound split at $|\zeta|=\sqrt t$.  Above this
scale,
\[
 |\widehat F(\zeta)|
 \le C\sqrt t\log^2(1/t)+\varepsilon_L=o(t^\kappa).
\]
By making the disorder threshold smaller, we obtain
$\widehat X(\zeta)\le(1+t^\kappa)p^{-2}$ on this range.
Below it, $t/|\zeta|^2\ge1$ and the rough upper bound is at most
$Ct/|\zeta|^2$.  Adding the nonnegative high-frequency term and using
\eqref{eq:initial-laplacian-symbol} proves the last line of
\eqref{eq:initial-comparisons} with its stated coefficient
$1+t^\kappa$.

All smallness conditions used above hold once $t_{n_0}$ is sufficiently
small.  Since $t_{n_0}\asymp4^{-n_0}$ uniformly in $m\ge4n_0$, one
choice of $\lambda_0(d)$ ensures them all. This completes the construction
of the initial state.
\end{proof}

Proposition~\ref{prop:initialization} supplies every initial comparison constant with fixed
slack.  Lemma~\ref{lem:dual-forms} also gives $X\asymp W^*$ and
$D\asymp W$ at this state.  The subsequent construction transports
these forms as sites are duplicated, updated, and eliminated; the
weighted cube estimate provides the corresponding local energy control.
